\documentclass[10pt,a4paper,final]{article} %draft

\usepackage[english, russian]{babel}

\usepackage{geometry}
%\usepackage[a4paper,left=2cm,right=2cm,top=2cm,bottom=2cm,bindingo ffset=0cm]{geometry}

\usepackage[T2A]{fontenc}
\usepackage[utf8]{inputenc}

\usepackage[utf8]{inputenc}

\usepackage[final]{pdfpages}

\usepackage{textcomp,enumitem}

\usepackage{amsmath,amsthm,amssymb}

\usepackage{longtable}
\usepackage{array}
\usepackage{adjustbox}




\usepackage{listings}
\usepackage{xcolor}
\usepackage{caption}

\definecolor{codegray}{rgb}{0.95,0.95,0.95}
\definecolor{codepurple}{rgb}{0.58,0,0.82}
\definecolor{codegreen}{rgb}{0,0.5,0} % зеленый цвет
\definecolor{codeblue}{rgb}{0,0,0.5} % синий цвет
\definecolor{codestring}{rgb}{0.64,0.08,0.08} % цвет строк


\lstset{
	language=C++,
	extendedchars=true,
	literate=
	{а}{{\selectfont\char224}}1
	{б}{{\selectfont\char225}}1
	{в}{{\selectfont\char226}}1
	{г}{{\selectfont\char227}}1
	{д}{{\selectfont\char228}}1
	{е}{{\selectfont\char229}}1
	{ж}{{\selectfont\char230}}1
	{з}{{\selectfont\char231}}1
	{и}{{\selectfont\char232}}1
	{й}{{\selectfont\char233}}1
	{к}{{\selectfont\char234}}1
	{л}{{\selectfont\char235}}1
	{м}{{\selectfont\char236}}1
	{н}{{\selectfont\char237}}1
	{о}{{\selectfont\char238}}1
	{п}{{\selectfont\char239}}1
	{р}{{\selectfont\char240}}1
	{с}{{\selectfont\char241}}1
	{т}{{\selectfont\char242}}1
	{у}{{\selectfont\char243}}1
	{ф}{{\selectfont\char244}}1
	{х}{{\selectfont\char245}}1
	{ц}{{\selectfont\char246}}1
	{ч}{{\selectfont\char247}}1
	{ш}{{\selectfont\char248}}1
	{щ}{{\selectfont\char249}}1
	{ъ}{{\selectfont\char250}}1
	{ы}{{\selectfont\char251}}1
	{ь}{{\selectfont\char252}}1
	{э}{{\selectfont\char253}}1
	{ю}{{\selectfont\char254}}1
	{я}{{\selectfont\char255}}1
	{А}{{\selectfont\char192}}1
	{Б}{{\selectfont\char193}}1
	{В}{{\selectfont\char194}}1
	{Г}{{\selectfont\char195}}1
	{Д}{{\selectfont\char196}}1
	{Е}{{\selectfont\char197}}1
	{Ж}{{\selectfont\char198}}1
	{З}{{\selectfont\char199}}1
	{И}{{\selectfont\char200}}1
	{Й}{{\selectfont\char201}}1
	{К}{{\selectfont\char202}}1
	{Л}{{\selectfont\char203}}1
	{М}{{\selectfont\char204}}1
	{Н}{{\selectfont\char205}}1
	{О}{{\selectfont\char206}}1
	{П}{{\selectfont\char207}}1
	{Р}{{\selectfont\char208}}1
	{С}{{\selectfont\char209}}1
	{Т}{{\selectfont\char210}}1
	{У}{{\selectfont\char211}}1
	{Ф}{{\selectfont\char212}}1
	{Х}{{\selectfont\char213}}1
	{Ц}{{\selectfont\char214}}1
	{Ч}{{\selectfont\char215}}1
	{Ш}{{\selectfont\char216}}1
	{Щ}{{\selectfont\char217}}1
	{Ъ}{{\selectfont\char218}}1
	{Ы}{{\selectfont\char219}}1
	{Ь}{{\selectfont\char220}}1
	{Э}{{\selectfont\char221}}1
	{Ю}{{\selectfont\char222}}1
	{Я}{{\selectfont\char223}}1,
	%	{|}{{\textbar}}1,
	%	{||}{{\textbar\textbar}}1
	%	{\&}{{\string&}}1
	%	{==}{{\string==}}1
	%	{\string\\}{{\string\\}}1,
	numbers=left,
	stepnumber=1,
	firstnumber=1,
	numberstyle=\tiny,
	basicstyle=\ttfamily\footnotesize,
	breakatwhitespace=false,
	breaklines=true,
	captionpos=b,
	keepspaces=true,
	showspaces=false,
	showstringspaces=false,
	showtabs=false,
	tabsize=2,
	frame=single,
	keywordstyle=\color{codeblue}, % ключевые слова
	commentstyle=\color{codegreen}, % комментарии
	stringstyle=\color{codestring}, % строки
	%identifierstyle=\color{green}, % стиль для переменных
	backgroundcolor=\color{codegray},
}
%\usepackage{fancyvrb}

%\usepackage{fancyhdr}

%\usepackage{upgreek}

%\usepackage{tipa}

%\usepackage{tikz}

%\usepackage{graphicx}

%\usepackage{pgfplots}

\usepackage{indentfirst}

%\usepackage{gensymb}

%\usepackage{amssymb}

%\usepackage{pdfpages}

\usepackage[unicode, pdftex, colorlinks, urlcolor=blue]{hyperref}

%\usepackage[T2A]{fontenc}


\usepackage{tabularx}
\usepackage{multirow}

\usepackage{booktabs} % для стильных линий таблицы



\usepackage{graphics}

\linespread{1.5}

\pagestyle{plain}

%\usepackage{listings}
%\usepackage{moreverb}

%\setlist[enumerate,itemize]{leftmargin=1.2cm} %отступ в перечислениях

\hypersetup{
	colorlinks=true,
	linkcolor=black,
	%allcolors=[RGB]{0 0 0}}
	}
	%\lstloadlanguages{ [LaTeX] TeX}

	\lstloadlanguages{ [LaTeX] TeX}



	\textheight=24cm
	\textwidth=16cm
	\oddsidemargin=0pt
	\topmargin=-1.5cm
	\parindent=24pt
	\parskip=0pt
	\tolerance=2000
	\flushbottom

	%\usepackage[font=scriptsize]{caption}
	\usepackage[labelsep=period]{caption}

	\begin{document}
\thispagestyle{empty}

\begin{center}
	{\Large МИНОБРНАУКИ РОССИИ}\\
	~\\
	{\large ФЕДЕРАЛЬНОЕ ГОСУДАРСТВЕННОЕ БЮДЖЕТНОЕ ОБРАЗОВАТЕЛЬНОЕ УЧРЕЖДЕНИЕ ВЫСШЕГО ПРОФЕССИОНАЛЬНОГО ОБРАЗОВАНИЯ}\\
	~\\
	{\Large \bf <<САНКТ-ПЕТЕРБУРГСКИЙ ПОЛИТЕХНИЧЕСКИЙ УНИВЕРСИТЕТ ПЕТРА ВЕЛИКОГО>>}\\
	~\\
	{\large Институт компьютерных наук и кибербезопасности}\\
	{\large Высшая школа технологий искусственного интеллекта}\\
	{\large Направление 02.03.01 Математика и компьютерные науки}\\
	~\\
	~\\
	~\\
	~\\
	{\Large \bf Отчет по лабораторной работе №1}\\
		{\Large \bf <<Создание лексического анализатора>>}\\
	{\Large по дисциплине <<Математическая логика и теория формальных языков>> }\\
	~\\
	~\\

	{\Large Вариант №44}\\
	%	{\Large \bf }\\

	~\\
	~\\
	~\\
	~\\
	~\\
	~\\
	~\\
	{\large Обучающийся: \underline{\hspace{3.5cm}} \qquad\qquad Гладков И.А.}\\
	~\\
	{\large Руководитель: \underline{\hspace{3.5cm}} \hspace{20mm} Востров А.В.}\\
	~\\
	~\\
\end{center}
\begin{flushright}

	«\underline{\hspace{1cm}}»\underline{\hspace{3cm}}20\underline{\hspace{0.7cm}}г.
\end{flushright}
~\\
\begin{center}
	{\large Санкт-Петербург, 2026}
\end{center}
\newpage

\tableofcontents

\newpage
\section* {Введение}
\addcontentsline{toc}{section}{Введение}



Целью проекта является разработка Telegram-бота для управления учебным расписанием, экзаменами, учебными задачами и напоминаниями. Бот должен позволять пользователю импортировать расписание из файла или внешнего источника, просматривать занятия и экзамены на сегодня, завтра и текущую неделю, вручную добавлять и изменять события, вести список учебных задач и получать рекомендации от LLM-советника по учебным приоритетам.

Система предназначена для студентов, которым необходимо централизованно хранить расписание занятий, экзамены, дедлайны и получать напоминания о важных учебных событиях.

\subsection{Основные функции}

\begin{enumerate}
    \item \textbf{Управление расписанием.}
    \begin{itemize}
        \item импорт расписания из файла CSV, iCal или внешнего API;
        \item просмотр расписания на сегодня, завтра и текущую неделю;
        \item ручное добавление занятий и экзаменов;
        \item редактирование занятий и экзаменов;
        \item удаление занятий и экзаменов.
    \end{itemize}

    \item \textbf{Управление учебными задачами.}
    \begin{itemize}
        \item добавление задач: упражнения, лабораторные, курсовые, дедлайны;
        \item привязка задач к учебным предметам;
        \item просмотр задач по предмету;
        \item просмотр задач по дедлайну: сегодня, завтра, неделя;
        \item отметка задачи как выполненной.
    \end{itemize}

    \item \textbf{Напоминания.}
    \begin{itemize}
        \item напоминания о предстоящих занятиях;
        \item напоминания о предстоящих экзаменах;
        \item напоминания о приближающихся дедлайнах задач.
    \end{itemize}

    \item \textbf{LLM-советник.}
    \begin{itemize}
        \item рекомендации по учебным приоритетам;
        \item генерация плана подготовки к экзамену;
        \item ответы на свободные вопросы пользователя по учебному планированию.
    \end{itemize}

    \item \textbf{Дополнительные функции.}
    \begin{itemize}
        \item интеграция с Google Calendar или Yandex Calendar;
        \item синхронизация расписания и задач с внешним календарем;
        \item учет времени обучения по предметам;
        \item LLM-анализ учебных привычек и рекомендации по улучшению.
    \end{itemize}
\end{enumerate}

\section{Общие требования}

\subsection{Управление версиями и совместная работа}

\begin{enumerate}[label=\arabic*.]
    \item Весь код проекта должен быть размещен в репозитории GitHub.
    \item Репозиторий должен содержать файл \cmd{README.md} с инструкцией по:
    \begin{itemize}
        \item настройке проекта локально;
        \item сборке и запуску приложения;
        \item запуску приложения с использованием Docker;
        \item настройке переменных окружения;
        \item подключению Telegram-бота;
        \item использованию основных команд бота;
        \item запуску тестов.
    \end{itemize}
    \item История коммитов должна отражать вклад всех участников команды.
    \item В репозитории должны отсутствовать секретные данные: токены Telegram, ключи LLM API, ключи календарных API, логины, пароли и строки подключения к базе данных.
\end{enumerate}

\subsection{Целевые артефакты}

\begin{enumerate}[label=\arabic*.]
    \item Приложение должно собираться в исполняемый \textbf{fat JAR}.
    \item Приложение должно быть развертываемо в виде Docker-контейнера.
    \item Для локального запуска должен быть предоставлен файл \cmd{docker-compose.yml}.
    \item Docker Compose должен запускать:
    \begin{itemize}
        \item приложение Telegram-бота;
        \item базу данных;
        \item при необходимости - дополнительные сервисы, например брокер сообщений или планировщик задач.
    \end{itemize}
    \item Docker-образ приложения должен быть опубликован в Docker Hub или другом согласованном контейнерном registry.
\end{enumerate}

\section{Функциональные требования}

\subsection{Аутентификация пользователей}

\begin{enumerate}[label=\arabic*.]
    \item Пользователь идентифицируется автоматически по Telegram \field{chat\_id} или \field{telegram\_id}.
    \item Дополнительная регистрация пользователя не требуется.
    \item При первом запуске бота создается запись пользователя в базе данных.
    \item Для каждого пользователя хранятся:
    \begin{itemize}
        \item Telegram ID;
        \item username, если он доступен;
        \item настройки часового пояса;
        \item настройки напоминаний;
        \item расписание;
        \item задачи;
        \item история учебных сессий, если реализована бонусная функция.
    \end{itemize}
\end{enumerate}

\subsection{Роли пользователей}

Система должна поддерживать две роли:

\begin{enumerate}
    \item \role{USER} - обычный пользователь Telegram-бота.
    \item \role{ADMIN} - администратор, имеющий доступ к служебным HTTP-эндпоинтам.
\end{enumerate}

\subsection{Авторизация администратора}

\begin{enumerate}[label=\arabic*.]
    \item Пользовательская часть работает через Telegram и не требует ручного ввода логина и пароля.
    \item Административная часть должна быть доступна через HTTP API.
    \item Для доступа к административным эндпоинтам используется API-ключ или Bearer-токен.
    \item API-ключ администратора должен выдаваться через эндпоинт \cmd{/auth} при передаче корректных данных администратора.
    \item Секретные данные администратора должны храниться в защищенном виде.
\end{enumerate}

\subsection{Служебные команды и эндпоинты}

\subsubsection{Healthcheck}

Команда или эндпоинт:

\begin{Code}
/healthcheck
GET /healthcheck
\end{Code}

Авторизация не требуется. Система должна возвращать:

\begin{itemize}
    \item состояние сервера;
    \item состояние подключения к базе данных;
    \item версию приложения;
    \item список авторов проекта.
\end{itemize}

Пример ответа:

\begin{Code}
Status: OK
Database: OK
Version: 1.0.0
Authors: Иванов И.И., Петров П.П., Сидоров С.С.
\end{Code}

\subsubsection{Получение списка пользователей}

HTTP-эндпоинт:

\begin{Code}
GET /users
\end{Code}

Доступен только пользователю с ролью \role{ADMIN}. Система должна возвращать список пользователей с их Telegram ID и базовой служебной информацией.

\section{Подробное описание функций}

\subsection{Запуск бота}

Команда:

\begin{Code}
/start
\end{Code}

Пользователь запускает Telegram-бота командой \cmd{/start}.

\textbf{Реакция системы:}

\begin{enumerate}
    \item Бот определяет пользователя по Telegram ID.
    \item Если пользователь новый, создается запись в базе данных.
    \item Бот отправляет приветственное сообщение и список доступных команд.
\end{enumerate}

Пример сообщения:

\begin{Code}
Привет, {username}! Я помогу тебе управлять расписанием занятий, экзаменами, учебными задачами и дедлайнами.

Доступные команды:
/menu - открыть меню
/today - расписание на сегодня
/tomorrow - расписание на завтра
/week - расписание на неделю
/import\_timetable - импорт расписания
/add\_lesson - добавить занятие
/add\_exam - добавить экзамен
/tasks - учебные задачи
/reminders - напоминания
/advisor - LLM-советник
/help - справка
\end{Code}

\subsection{Меню}

Команда:

\begin{Code}
/menu
\end{Code}

Бот выводит основное меню с кнопками:

\begin{enumerate}
    \item Расписание.
    \item Задачи.
    \item Напоминания.
    \item LLM-советник.
    \item Импорт расписания.
    \item Настройки.
    \item Справка.
\end{enumerate}

Команды меню могут быть реализованы как Telegram-команды или inline-кнопки.

\section{Расписание занятий и экзаменов}

\subsection{Импорт расписания}

Команда:

\begin{Code}
/import\_timetable
\end{Code}

Пользователь отправляет команду \cmd{/import\_timetable}, после чего загружает файл с расписанием или указывает ссылку на внешний источник.

\subsubsection{Поддерживаемые источники}

Система должна поддерживать хотя бы один из следующих способов импорта:

\begin{enumerate}
    \item CSV-файл.
    \item iCal-файл.
    \item Внешний API учебного заведения.
\end{enumerate}

Желательно поддержать CSV и iCal как базовые форматы.

\subsubsection{Формат CSV}

CSV-файл должен содержать следующие поля:

\begin{Code}
type,subject,title,start_datetime,end_datetime,location,teacher,comment
\end{Code}

Описание полей:

\begin{itemize}
    \item \field{type} - тип события: \field{lesson} или \field{exam};
    \item \field{subject} - название предмета;
    \item \field{title} - название занятия или экзамена;
    \item \field{start\_datetime} - дата и время начала;
    \item \field{end\_datetime} - дата и время окончания;
    \item \field{location} - аудитория или ссылка на онлайн-занятие;
    \item \field{teacher} - преподаватель;
    \item \field{comment} - дополнительная информация.
\end{itemize}

Пример строки для занятия:

\begin{Code}
lesson,Математический анализ,Лекция,12.05.2026 10:00,12.05.2026 11:30,ауд. 305,Иванов И.И.,
\end{Code}

Пример строки для экзамена:

\begin{Code}
exam,Алгоритмы и структуры данных,Экзамен,20.06.2026 09:00,20.06.2026 12:00,ауд. 101,Петров П.П.,Взять зачётку
\end{Code}

\subsubsection{Корректный ввод}

Файл считается корректным, если:

\begin{itemize}
    \item имеет поддерживаемый формат;
    \item содержит обязательные поля;
    \item даты и время указаны в поддерживаемом формате;
    \item время окончания события позже времени начала;
    \item тип события равен \field{lesson} или \field{exam}.
\end{itemize}

\subsubsection{Некорректный ввод}

Файл считается некорректным, если:

\begin{itemize}
    \item формат файла не поддерживается;
    \item отсутствуют обязательные поля;
    \item дата или время не распознаются;
    \item событие содержит некорректный тип;
    \item время окончания раньше времени начала;
    \item файл пустой.
\end{itemize}

\subsubsection{Реакция системы}

Если импорт успешен:

\begin{Code}
Расписание успешно импортировано.
Добавлено событий: {count}.
Занятий: {lessons_count}.
Экзаменов: {exams_count}.
\end{Code}

Если часть записей некорректна:

\begin{Code}
Импорт выполнен частично.
Добавлено событий: {count}.
Пропущено строк: {skipped_count}.
Причины ошибок доступны в отчёте.
\end{Code}

Если импорт невозможен:

\begin{Code}
Не удалось импортировать расписание. Проверьте формат файла и повторите попытку.
\end{Code}

\subsection{Просмотр расписания на сегодня}

Команда:

\begin{Code}
/today
\end{Code}

Бот выводит список занятий и экзаменов пользователя на текущую дату. События должны быть отсортированы по времени начала.

Формат ответа:

\begin{Code}
Расписание на сегодня, {date}:

1. 10:00-11:30
Математический анализ
Тип: Лекция
Место: ауд. 305
Преподаватель: Иванов И.И.

2. 13:00-14:30
Алгоритмы и структуры данных
Тип: Практика
Место: ауд. 204
Преподаватель: Петров П.П.
\end{Code}

Если событий нет:

\begin{Code}
На сегодня занятий и экзаменов нет.
\end{Code}

\subsection{Просмотр расписания на завтра}

Команда:

\begin{Code}
/tomorrow
\end{Code}

Бот выводит список занятий и экзаменов на следующий календарный день. Если событий нет, бот отправляет сообщение:

\begin{Code}
На завтра занятий и экзаменов нет.
\end{Code}

\subsection{Просмотр расписания на неделю}

Команда:

\begin{Code}
/week
\end{Code}

Бот выводит расписание на текущую неделю. Текущая неделя считается от понедельника до воскресенья, если иное не задано в настройках пользователя. События группируются по датам.

Пример ответа:

\begin{Code}
Расписание на неделю:

Понедельник, 11.05.2026
10:00-11:30 - Математический анализ, лекция, ауд. 305

Среда, 13.05.2026
13:00-14:30 - Базы данных, практика, ауд. 212

Пятница, 15.05.2026
09:00-10:30 - Английский язык, семинар, ауд. 110
\end{Code}

Если событий нет:

\begin{Code}
На этой неделе занятий и экзаменов нет.
\end{Code}

\subsection{Добавление занятия}

Команда:

\begin{Code}
/add_lesson
\end{Code}

Пользователь вводит команду \cmd{/add\_lesson} и по шагам заполняет данные занятия.

\textbf{Обязательные поля:}

\begin{itemize}
    \item предмет;
    \item тип занятия;
    \item дата;

\end{itemize}

\textbf{Необязательные поля:}

\begin{itemize}
    \item время начала;
    \item время окончания.
    \item аудитория или ссылка;
    \item преподаватель;
    \item комментарий.
\end{itemize}

Пример корректного ввода:

\begin{Code}
Предмет: Математический анализ
Тип: Лекция
Дата: 12.05.2026
Время начала: 10:00
Время окончания: 11:30
Место: ауд. 305
Преподаватель: Иванов И.И.
\end{Code}

При успешном добавлении:

\begin{Code}
Занятие успешно добавлено в расписание.

Математический анализ
Лекция
12.05.2026, 10:00-11:30
Место: ауд. 305
\end{Code}

При некорректном вводе:

\begin{Code}
Не удалось добавить занятие. Проверьте дату, время и обязательные поля.
\end{Code}

\subsection{Добавление экзамена}

Команда:

\begin{Code}
/add_exam
\end{Code}

\textbf{Обязательные поля:}

\begin{itemize}
    \item предмет;
    \item дата;

\end{itemize}

\textbf{Необязательные поля:}

\begin{itemize}
\item время начала;
    \item место проведения.
    \item время окончания;
    \item преподаватель или экзаменатор;
    \item комментарий;
    \item список тем;
    \item ссылка на материалы для подготовки.
\end{itemize}

При успешном добавлении:

\begin{Code}
Экзамен успешно добавлен.

Предмет: Алгоритмы и структуры данных
Дата: 20.06.2026
Время: 09:00
Место: ауд. 101
\end{Code}

\subsection{Редактирование события}

Команда:

\begin{Code}
/edit_event
\end{Code}

или отдельные команды:

\begin{Code}
/edit_lesson
/edit_exam
\end{Code}

Пользователь выбирает событие из списка или вводит его идентификатор.

\textbf{Реакция системы:}

\begin{enumerate}
    \item Бот показывает текущие данные события.
    \item Бот предлагает выбрать поле для изменения:
    \begin{itemize}
        \item предмет;
        \item дата;
        \item время;
        \item место;
        \item преподаватель;
        \item комментарий.
    \end{itemize}
    \item Пользователь вводит новое значение.
    \item Система валидирует данные.
    \item Событие обновляется.
\end{enumerate}

При успехе:

\begin{Code}
Событие успешно обновлено.
\end{Code}

Если событие не найдено:

\begin{Code}
Событие не найдено. Проверьте выбранный идентификатор.
\end{Code}

\subsection{Удаление события}

Команда:

\begin{Code}
/delete_event
\end{Code}

или отдельные команды:

\begin{Code}
/delete_lesson
/delete_exam
\end{Code}

\textbf{Реакция системы:}

\begin{enumerate}
    \item Бот показывает ближайшие события пользователя.
    \item Пользователь выбирает событие.
    \item Бот запрашивает подтверждение удаления.
    \item После подтверждения событие удаляется.
\end{enumerate}

При успехе:

\begin{Code}
Событие удалено из расписания.
\end{Code}

При отмене:

\begin{Code}
Удаление отменено.
\end{Code}

\section{Управление учебными задачами}

\subsection{Меню задач}

Команда:

\begin{Code}
/tasks
\end{Code}

Бот показывает меню задач:

\begin{Code}
Управление задачами:

/add_task - добавить задачу
/tasks_today - задачи на сегодня
/tasks_tomorrow - задачи на завтра
/tasks_week - задачи на неделю
/tasks_by_subject - задачи по предмету
/done_task - отметить задачу выполненной
\end{Code}

\subsection{Добавление задачи}

Команда:

\begin{Code}
/add_task
\end{Code}

\textbf{Обязательные поля:}

\begin{itemize}
    \item название задачи;
    \item предмет;

\end{itemize}

\textbf{Необязательные поля:}

\begin{itemize}
\item дедлайн.
    \item описание;
    \item тип задачи;
    \item приоритет;
    \item примерное время выполнения.
\end{itemize}

\textbf{Типы задач:}

\begin{itemize}
    \item упражнение;
    \item лабораторная работа;
    \item домашнее задание;
    \item курсовая работа;
    \item подготовка к экзамену;
    \item чтение материалов;
    \item другое.
\end{itemize}

Пример корректного ввода:

\begin{Code}
Название: Решить задачи 1-10
Предмет: Математический анализ
Тип: Домашнее задание
Дедлайн: 15.05.2026 23:59
Приоритет: высокий
Описание: задачи из методички, страница 42
\end{Code}

При успехе:

\begin{Code}
Задача успешно добавлена.

Решить задачи 1-10
Предмет: Математический анализ
Дедлайн: 15.05.2026 23:59
Приоритет: высокий
\end{Code}

При ошибке:

\begin{Code}
Не удалось добавить задачу. Проверьте название, предмет и дедлайн.
\end{Code}

\subsection{Просмотр задач на сегодня}

Команда:

\begin{Code}
/tasks_today
\end{Code}

Бот выводит все невыполненные задачи, дедлайн которых наступает сегодня.

Формат ответа:

\begin{Code}
Задачи на сегодня:

1. Математический анализ
Решить задачи 1-10
Дедлайн: 23:59
Приоритет: высокий

2. Базы данных
Подготовить ER-диаграмму
Дедлайн: 18:00
Приоритет: средний
\end{Code}

Если задач нет:

\begin{Code}
На сегодня задач нет.
\end{Code}

\subsection{Просмотр задач на завтра}

Команда:

\begin{Code}
/tasks_tomorrow
\end{Code}

Бот выводит все невыполненные задачи с дедлайном на завтра. Если задач нет:

\begin{Code}
На завтра задач нет.
\end{Code}

\subsection{Просмотр задач на неделю}

Команда:

\begin{Code}
/tasks_week
\end{Code}

Бот выводит задачи, дедлайн которых находится в пределах текущей недели. Задачи должны быть отсортированы по дедлайну и приоритету. Если задач нет:

\begin{Code}
На этой неделе задач нет.
\end{Code}

\subsection{Просмотр задач по предмету}

Команда:

\begin{Code}
/tasks_by_subject
\end{Code}

Пользователь вводит название предмета или выбирает его из списка. Бот выводит невыполненные задачи по выбранному предмету.

Если предмет не найден:

\begin{Code}
Предмет не найден. Проверьте название или добавьте задачу/занятие по этому предмету.
\end{Code}

Если задач по предмету нет:

\begin{Code}
По предмету {subject} активных задач нет.
\end{Code}

\subsection{Отметка задачи как выполненной}

Команда:

\begin{Code}
/done_task
\end{Code}

\textbf{Реакция системы:}

\begin{enumerate}
    \item Бот показывает список активных задач.
    \item Пользователь выбирает задачу.
    \item Система меняет статус задачи на \field{DONE}.
\end{enumerate}

При успехе:

\begin{Code}
Задача отмечена как выполненная.
\end{Code}

Если задача не найдена:

\begin{Code}
Задача не найдена или уже выполнена.
\end{Code}

\section{Напоминания}

\subsection{Общие требования к напоминаниям}

\begin{enumerate}[label=\arabic*.]
    \item Система должна отправлять пользователю Telegram-уведомления о предстоящих занятиях, экзаменах и дедлайнах.
    \item Пользователь должен иметь возможность включать и отключать напоминания.
    \item Напоминания должны учитывать часовой пояс пользователя.
    \item По умолчанию должны быть включены следующие напоминания:
    \begin{itemize}
        \item за 30 минут до занятия;
        \item за 1 день до экзамена;
        \item за 2 часа до дедлайна задачи.
    \end{itemize}
\end{enumerate}

Значения по умолчанию могут быть изменены в настройках.

\subsection{Просмотр настроек напоминаний}

Команда:

\begin{Code}
/reminders
\end{Code}

Бот выводит текущие настройки:

\begin{Code}
Настройки напоминаний:

Занятия: за 30 минут
Экзамены: за 1 день
Задачи: за 2 часа
Статус: включены

Команды:
/set_reminder - изменить настройки
/disable_reminders - отключить напоминания
/enable_reminders - включить напоминания
\end{Code}

\subsection{Изменение напоминаний}

Команда:

\begin{Code}
/set_reminder
\end{Code}

Бот предлагает выбрать тип напоминания:

\begin{enumerate}
    \item Занятие.
    \item Экзамен.
    \item Задача.
\end{enumerate}

После выбора пользователь вводит время до события.

Пример корректного ввода:

\begin{Code}
30m
2h
1d
\end{Code}

Где:

\begin{itemize}
    \item \field{m} - минуты;
    \item \field{h} - часы;
    \item \field{d} - дни.
\end{itemize}

При успехе:

\begin{Code}
Настройки напоминания обновлены.
\end{Code}

При ошибке:

\begin{Code}
Некорректный формат времени. Используйте формат 30m, 2h или 1d.
\end{Code}

\subsection{Напоминание о занятии}

Система должна отправлять сообщение:

\begin{Code}
Напоминание о занятии:

Через 30 минут начнётся занятие:
Математический анализ
Тип: Лекция
Время: 10:00-11:30
Место: ауд. 305
\end{Code}

\subsection{Напоминание об экзамене}

Система должна отправлять сообщение:

\begin{Code}
Напоминание об экзамене:

Завтра экзамен:
Алгоритмы и структуры данных
Дата: 20.06.2026
Время: 09:00
Место: ауд. 101
\end{Code}

\subsection{Напоминание о дедлайне}

Система должна отправлять сообщение:

\begin{Code}
Напоминание о дедлайне:

До дедлайна осталось 2 часа.
Задача: Подготовить ER-диаграмму
Предмет: Базы данных
Дедлайн: 18:00
\end{Code}

\section{LLM-советник}

\subsection{Общие требования}

\begin{enumerate}[label=\arabic*.]
    \item Бот должен использовать внешний LLM API для генерации учебных рекомендаций.
    \item LLM-советник должен использовать данные пользователя:
    \begin{itemize}
        \item расписание занятий;
        \item даты экзаменов;
        \item учебные задачи;
        \item дедлайны;
        \item приоритеты задач;
        \item историю учебных сессий, если реализована бонусная функция.
    \end{itemize}
    \item Ответы LLM должны носить рекомендательный характер.
    \item Если данных недостаточно, бот должен сообщить пользователю, какие данные нужно добавить.
\end{enumerate}

Пример сообщения при недостатке данных:

\begin{Code}
Недостаточно данных для рекомендации. Добавьте задачи или экзамены, чтобы я мог оценить приоритеты.
\end{Code}

\subsection{Меню LLM-советника}

Команда:

\begin{Code}
/advisor
\end{Code}

Бот выводит меню:

\begin{Code}
LLM-советник:

/priority - что учить в первую очередь
/exam_plan - план подготовки к экзамену
/ask_advisor - задать вопрос советнику
\end{Code}

\subsection{Рекомендация по учебным приоритетам}

Команда:

\begin{Code}
/priority
\end{Code}

Система анализирует:

\begin{itemize}
    \item ближайшие дедлайны;
    \item даты экзаменов;
    \item количество невыполненных задач;
    \item приоритет задач;
    \item свободные временные окна в расписании.
\end{itemize}

Пример ответа:

\begin{Code}
Рекомендация по приоритетам:

1. Базы данных - высокий приоритет
Причина: дедлайн по ER-диаграмме сегодня в 18:00.

2. Алгоритмы и структуры данных - высокий приоритет
Причина: экзамен через 5 дней, есть 3 незавершённые задачи.

3. Математический анализ - средний приоритет
Причина: ближайший дедлайн через 4 дня.
\end{Code}

\subsection{Генерация плана подготовки к экзамену}

Команда:

\begin{Code}
/exam_plan
\end{Code}

Пользователь выбирает экзамен из списка или вводит название предмета. Бот формирует план подготовки на основе:

\begin{itemize}
    \item даты экзамена;
    \item оставшегося количества дней;
    \item списка задач по предмету;
    \item свободных окон в расписании;
    \item указанной пользователем сложности предмета, если такая информация есть.
\end{itemize}

Пример ответа:

\begin{Code}
План подготовки к экзамену по предмету "Алгоритмы и структуры данных":

День 1:
Повторить сортировки и асимптотику.
Время: 1.5 часа.

День 2:
Разобрать деревья поиска и хеш-таблицы.
Время: 2 часа.

День 3:
Решить тренировочные задачи.
Время: 2 часа.

День 4:
Повторить конспекты и слабые темы.
Время: 1 час.

День 5:
Лёгкое повторение, без перегрузки.
Время: 1 час.
\end{Code}

Если экзамен не найден:

\begin{Code}
Экзамен по указанному предмету не найден. Добавьте экзамен через /add_exam.
\end{Code}

\subsection{Свободный вопрос LLM-советнику}

Команда:

\begin{Code}
/ask_advisor
\end{Code}

Пользователь задает вопрос в свободной форме. Примеры:

\begin{Code}
Как распределить подготовку к двум экзаменам?
Что лучше учить сегодня?
Как не забыть про дедлайны на этой неделе?
\end{Code}

\textbf{Реакция системы:}

\begin{enumerate}
    \item Бот получает вопрос пользователя.
    \item Бот добавляет в контекст расписание, задачи и экзамены пользователя.
    \item Бот отправляет запрос к LLM.
    \item Бот возвращает ответ пользователю.
\end{enumerate}

Если LLM API недоступен:

\begin{Code}
LLM-советник временно недоступен. Попробуйте повторить запрос позже.
\end{Code}

\section{Дополнительные функции}

\subsection{Интеграция с календарями}

В качестве бонусной функции система может поддерживать интеграцию с:

\begin{itemize}
    \item Google Calendar;
    \item Yandex Calendar;
    \item другим календарным API, доступным в учебном заведении.
\end{itemize}

Система должна позволять:

\begin{itemize}
    \item экспортировать занятия в календарь;
    \item экспортировать экзамены в календарь;
    \item экспортировать дедлайны задач в календарь;
    \item синхронизировать изменения расписания;
    \item отключать синхронизацию.
\end{itemize}

Команды:

\begin{Code}
/calendar_sync
/calendar_export
/calendar_disable
\end{Code}

При успешной синхронизации:

\begin{Code}
Синхронизация с календарём выполнена.
Добавлено событий: {count}.
\end{Code}

При ошибке авторизации:

\begin{Code}
Не удалось подключиться к календарю. Проверьте доступы и повторите попытку.
\end{Code}

\subsection{Учет учебного времени}

Команды:

\begin{Code}
/study_start
/study_stop
/study_stats
\end{Code}

Пользователь может запускать и завершать учебную сессию по конкретному предмету.

\textbf{Сценарий:}

\begin{enumerate}
    \item Пользователь вводит команду \cmd{/study\_start}.
    \item Бот запрашивает предмет.
    \item Пользователь выбирает предмет.
    \item Бот фиксирует начало учебной сессии.
\end{enumerate}

Пример ответа:

\begin{Code}
Учебная сессия начата.
Предмет: Базы данных
Начало: 16:00
\end{Code}

Пользователь завершает сессию командой \cmd{/study\_stop}. Бот сохраняет длительность.

Пример ответа:

\begin{Code}
Учебная сессия завершена.
Предмет: Базы данных
Длительность: 1 час 20 минут
\end{Code}

Команда \cmd{/study\_stats} выводит:

\begin{itemize}
    \item общее время обучения за день;
    \item общее время обучения за неделю;
    \item распределение времени по предметам.
\end{itemize}

\subsection{LLM-анализ учебных привычек}

Команда:

\begin{Code}
/habits_analysis
\end{Code}

Если реализован учет учебного времени, бот может анализировать учебные привычки пользователя. Система передает в LLM:

\begin{itemize}
    \item историю учебных сессий;
    \item задачи;
    \item дедлайны;
    \item расписание;
    \item экзамены.
\end{itemize}

Пример ответа:

\begin{Code}
Анализ учебных привычек:

1. Ты чаще занимаешься вечером, но дедлайны часто закрываются в последний день.
2. Рекомендуется начинать подготовку к крупным задачам минимум за 3 дня.
3. На этой неделе больше всего времени ушло на Базы данных, но экзамен по Алгоритмам ближе.
4. Совет: выдели завтра 1.5 часа на Алгоритмы и 1 час на Математический анализ.
\end{Code}

\section{Ограничения и лимиты}

\begin{enumerate}[label=\arabic*.]
    \item Один пользователь не должен иметь более одного активного процесса импорта расписания одновременно.
    \item Максимальный размер импортируемого файла - 5 МБ.
    \item Максимальное количество событий расписания на одного пользователя - 5000.
    \item Максимальное количество активных задач на одного пользователя - 1000.
    \item Максимальная длина названия предмета - 100 символов.
    \item Максимальная длина описания задачи - 1000 символов.
    \item Максимальная длина вопроса к LLM-советнику - 2000 символов.
    \item Бот должен поддерживать не менее 1000 активных пользователей.
    \item Расписание на неделю должно выводиться не более чем за 2 секунды без учета задержек внешних API.
    \item При недоступности внешнего календарного API система должна сохранять данные локально и сообщать пользователю о проблеме синхронизации.
\end{enumerate}

\section{Нефункциональные требования}

\subsection{Производительность}

\begin{enumerate}[label=\arabic*.]
    \item Максимальное время ответа на простые команды, такие как \cmd{/menu}, \cmd{/today}, \cmd{/tasks\_today}, не должно превышать 500 мс без учета задержек внешних API.
    \item Импорт расписания должен выполняться не дольше 10 секунд для файла до 1000 событий.
    \item Запрос к LLM должен иметь таймаут не более 30 секунд.
    \item При долгих операциях бот должен сообщать пользователю, что запрос принят в обработку.
\end{enumerate}

\subsection{Надежность}

\begin{enumerate}[label=\arabic*.]
    \item При ошибках внешнего API система должна выполнять повторные попытки.
    \item Количество повторных попыток - не более 3.
    \item Интервал между повторными попытками - 1 секунда для обычных API и до 30 секунд для LLM API.
    \item Если внешний API недоступен после всех попыток, пользователь должен получить понятное сообщение об ошибке.
    \item Все данные пользователя должны сохраняться в базе данных.
    \item Система не должна терять уже сохраненное расписание при неудачном импорте нового файла.
\end{enumerate}

\subsection{Валидация ввода}

Система должна проверять:

\begin{itemize}
    \item корректность дат;
    \item корректность времени;
    \item наличие обязательных полей;
    \item корректность формата файла;
    \item корректность типа события;
    \item длину пользовательского ввода;
    \item принадлежность задачи существующему или новому предмету;
    \item корректность команд Telegram.
\end{itemize}

При некорректном вводе бот должен объяснять ошибку и предлагать повторить ввод.

\subsection{Безопасность}

\begin{enumerate}[label=\arabic*.]
    \item Токен Telegram-бота должен храниться в переменных окружения.
    \item Ключи LLM API должны храниться в переменных окружения.
    \item Ключи календарных API должны храниться в переменных окружения или защищенном хранилище.
    \item Данные подключения к базе данных не должны храниться в коде.
    \item Административные эндпоинты должны быть защищены авторизацией.
    \item Пользователь не должен иметь доступ к расписанию, задачам и учебным данным других пользователей.
    \item Все входные данные должны проходить валидацию для защиты от инъекций и некорректных запросов.
\end{enumerate}

\subsection{Логирование}

Система должна логировать:

\begin{itemize}
    \item запуск приложения;
    \item ошибки Telegram API;
    \item ошибки базы данных;
    \item ошибки внешних API;
    \item ошибки импорта расписания;
    \item ошибки LLM-запросов;
    \item административные запросы.
\end{itemize}

В логах не должны сохраняться:

\begin{itemize}
    \item токены;
    \item пароли;
    \item API-ключи;
    \item персональные данные в избыточном объеме.
\end{itemize}

\subsection{Обработка ошибок}

Все ошибки должны сопровождаться понятным сообщением для пользователя.

Примеры:

\begin{Code}
Файл не удалось обработать. Проверьте формат и повторите попытку.
\end{Code}

\begin{Code}
Событие не найдено. Возможно, оно было удалено ранее.
\end{Code}

\begin{Code}
LLM-советник временно недоступен. Попробуйте позже.
\end{Code}

\begin{Code}
Не удалось подключиться к календарю. Проверьте настройки интеграции.
\end{Code}

\section{Требования к технологическому стеку}

Рекомендуемый стек проекта:

\begin{itemize}
    \item язык программирования: Java;
    \item фреймворк: Spring;
    \item Telegram API: для реализации Telegram-бота;
    \item база данных: PostgreSQL или MongoDB;
    \item инструмент сборки: Maven или Gradle;
    \item контейнеризация: Docker;
    \item локальный запуск: Docker Compose;
    \item тестирование: JUnit;
    \item документация API: Spring REST Docs или Swagger/OpenAPI;
    \item внешний LLM API: любой доступный LLM-провайдер;
    \item календарные интеграции: Google Calendar API, Yandex Calendar API или API учебного заведения.
\end{itemize}

Конкретные версии Java, Spring и базы данных следует согласовать с требованиями курса или преподавателя.

\section{Требования к документации}

Документация проекта должна включать:

\begin{enumerate}
    \item \textbf{README.md.}
    \begin{itemize}
        \item описание проекта;
        \item список функций;
        \item список команд Telegram-бота;
        \item инструкцию по локальному запуску;
        \item инструкцию по запуску через Docker;
        \item описание переменных окружения;
        \item инструкцию по запуску тестов.
    \end{itemize}

    \item \textbf{Документация API.}
    \begin{itemize}
        \item описание эндпоинтов \cmd{/healthcheck}, \cmd{/auth}, \cmd{/users};
        \item описание формата запросов и ответов;
        \item описание ошибок.
    \end{itemize}

    \item \textbf{Документация по формату импорта.}
    \begin{itemize}
        \item пример CSV-файла;
        \item описание обязательных полей;
        \item описание поддерживаемых форматов даты и времени;
        \item пример iCal-файла или ссылка на поддерживаемый формат.
    \end{itemize}

    \item \textbf{Документация по архитектуре.}
    \begin{itemize}
        \item схема компонентов;
        \item описание модулей системы;
        \item схема взаимодействия с Telegram API;
        \item схема взаимодействия с LLM API;
        \item схема взаимодействия с календарным API, если реализовано.
    \end{itemize}
\end{enumerate}

\section{Руководство по сдаче проекта}

\subsection{Требования}

Необходимо предоставить документ с требованиями, включающий:

\begin{itemize}
    \item цель проекта;
    \item описание основной функциональности;
    \item функциональные требования;
    \item нефункциональные требования;
    \item ограничения;
    \item используемый технологический стек;
    \item глоссарий.
\end{itemize}

\subsection{Архитектура}

Необходимо предоставить документ по архитектуре, включающий:

\begin{itemize}
    \item высокоуровневое описание системы;
    \item компонентную диаграмму;
    \item описание базы данных;
    \item описание взаимодействия с Telegram API;
    \item описание взаимодействия с LLM API;
    \item описание работы напоминаний;
    \item описание развертывания приложения.
\end{itemize}

\subsection{Полный проект}

Необходимо предоставить:

\begin{itemize}
    \item ссылку на GitHub-репозиторий;
    \item ссылку на Docker Hub с опубликованным образом;
    \item имя или ссылку на Telegram-бота;
    \item инструкцию по тестированию;
    \item примеры файлов для импорта расписания.
\end{itemize}

\section{Сценарии использования}

\subsection{Импорт расписания из файла}

\textbf{Действия пользователя:}

\begin{enumerate}
    \item Пользователь вводит команду \cmd{/import\_timetable}.
    \item Пользователь отправляет CSV или iCal-файл.
\end{enumerate}

\textbf{Реакция системы:}

\begin{enumerate}
    \item Бот принимает файл.
    \item Проверяет формат.
    \item Валидирует строки расписания.
    \item Сохраняет корректные события в базу данных.
    \item Возвращает отчет об импорте.
\end{enumerate}

\textbf{Результат:} пользователь получает сообщение об успешном или частичном импорте расписания.

\subsection{Просмотр расписания на сегодня}

\textbf{Действия пользователя:}

\begin{enumerate}
    \item Пользователь вводит команду \cmd{/today}.
\end{enumerate}

\textbf{Реакция системы:}

\begin{enumerate}
    \item Бот определяет текущую дату с учетом часового пояса пользователя.
    \item Загружает занятия и экзамены на текущую дату.
    \item Сортирует события по времени.
    \item Отправляет пользователю список событий.
\end{enumerate}

\textbf{Результат:} пользователь видит расписание на сегодня.

\subsection{Добавление задачи}

\textbf{Действия пользователя:}

\begin{enumerate}
    \item Пользователь вводит команду \cmd{/add\_task}.
    \item Пользователь вводит название задачи.
    \item Пользователь выбирает или вводит предмет.
    \item Пользователь указывает дедлайн.
    \item Пользователь при необходимости указывает приоритет и описание.
\end{enumerate}

\textbf{Реакция системы:}

\begin{enumerate}
    \item Бот проверяет корректность данных.
    \item Сохраняет задачу в базе данных.
    \item Подтверждает добавление задачи.
\end{enumerate}

\textbf{Результат:} задача появляется в списке задач пользователя.

\subsection{Получение рекомендации по приоритетам}

\textbf{Действия пользователя:}

\begin{enumerate}
    \item Пользователь вводит команду \cmd{/priority}.
\end{enumerate}

\textbf{Реакция системы:}

\begin{enumerate}
    \item Бот собирает данные о задачах, дедлайнах, экзаменах и расписании.
    \item Передает структурированный контекст в LLM.
    \item Получает рекомендацию.
    \item Отправляет пользователю список приоритетов.
\end{enumerate}

\textbf{Результат:} пользователь получает рекомендацию, какие предметы и задачи нужно выполнить в первую очередь.

\subsection{Получение напоминания о дедлайне}

\textbf{Условие:} у пользователя есть активная задача с дедлайном.

\textbf{Реакция системы:}

\begin{enumerate}
    \item Планировщик проверяет ближайшие дедлайны.
    \item Если дедлайн приближается, система отправляет сообщение в Telegram.
    \item Пользователь получает напоминание.
\end{enumerate}

\textbf{Результат:} пользователь заранее узнает о приближающемся дедлайне.

\section{Глоссарий}

\begin{longtable}{p{0.30\textwidth}p{0.64\textwidth}}
\toprule
\textbf{Термин} & \textbf{Определение} \\
\midrule
\endhead
Telegram-бот & Автоматизированная программа, работающая внутри Telegram и обрабатывающая команды пользователя. \\
Telegram ID & Уникальный идентификатор пользователя в Telegram. \\
Расписание & Набор занятий и экзаменов пользователя с датой, временем и дополнительными параметрами. \\
Занятие & Учебное событие: лекция, практика, лабораторная, семинар или другое занятие. \\
Экзамен & Контрольное учебное событие по предмету, имеющее дату, время и место проведения. \\
Предмет & Учебная дисциплина, к которой относятся занятия, экзамены и задачи. \\
Задача & Учебное дело пользователя: домашняя работа, упражнение, курсовая, лабораторная или другой дедлайн. \\
Дедлайн & Дата и время, к которым задача должна быть выполнена. \\
Напоминание & Сообщение от бота, отправляемое до наступления занятия, экзамена или дедлайна. \\
LLM & Большая языковая модель, используемая для генерации учебных рекомендаций. \\
LLM-советник & Функция бота, которая помогает пользователю определить приоритеты и составить план подготовки. \\
CSV & Текстовый формат хранения табличных данных, используемый для импорта расписания. \\
iCal & Формат календарных событий, используемый для обмена расписанием между календарными приложениями. \\
Google Calendar API & API для интеграции с календарем Google. \\
Yandex Calendar API & API для интеграции с календарем Яндекса. \\
Docker & Платформа для контейнеризации приложения. \\
Docker Compose & Инструмент для запуска нескольких связанных Docker-контейнеров. \\
Fat JAR & Исполняемый JAR-файл, содержащий приложение и необходимые зависимости. \\
База данных & Система хранения данных о пользователях, расписании, задачах и напоминаниях. \\
Healthcheck & Проверка работоспособности приложения и его зависимостей. \\
API-ключ & Секретная строка, используемая для доступа к административным эндпоинтам. \\
ADMIN & Роль администратора системы. \\
USER & Стандартная роль пользователя Telegram-бота. \\
Учебная сессия & Период времени, в течение которого пользователь занимался конкретным предметом. \\
Синхронизация & Обмен событиями между ботом и внешним календарем. \\
\bottomrule
\end{longtable}

\end{document}


\end{document}